The gap analysis points to a concrete research and standardization agenda. We organize it
by priority.

\textbf{Action-level evaluation for agentic systems.} The agentic measurement gap
(Section~\ref{sec:gaps}) is the most urgent. The field needs evaluation harnesses that
instrument an agent's \emph{trajectory}---tool calls, delegations, plan steps---and
compute disparity, safety, and reliability over actions rather than tokens. Such
harnesses would simultaneously serve EU Article~12 logging, Article~14 oversight, and
NIST \textsc{Manage} monitoring, turning a compliance burden into a measurement
substrate.

\textbf{Standardized conformity evaluations and harmonised benchmarks.} At present,
mapping a method to a requirement still requires expert interpretation. The natural next
step is \emph{harmonised standards}---the CEN--CENELEC deliverables anticipated by the EU
AI Act, and NIST/ISO test methods---that fix, for each obligation, a reference evaluation
protocol and an acceptance criterion. A living, versioned benchmark registry tied to
specific clauses would reduce the interpretive latitude that our matrix currently
documents.

\textbf{Evaluation of process and oversight obligations.} The under-served, non-process
gaps---oversight effectiveness, log adequacy, impact-assessment quality---call for new
evaluation \emph{methods}, not merely new documentation. Human-subjects protocols for
appropriate reliance, log-reconstruction audits, and quantitative impact-assessment
instruments would convert today's documentary evidence into measurable evidence.

\textbf{Third-party audit and assurance ecosystems.} Conformity assessment ultimately
depends on credible, independent evaluation. Recent work on rigorous third-party auditing
of frontier developers~\cite{auditing2024frontier} and scalable multi-policy compliance
evaluation~\cite{compliance2025pasta} points toward an assurance ecosystem; in the US,
this aligns with a NIST-anchored, voluntary-conformity model of national importance for
AI safety and governance.

\textbf{Automated and continuous evidence generation.} As models and regulations both
evolve, one-shot evaluation is insufficient. Continuous-assurance pipelines that
re-generate conformity evidence on each model or data update---and that flag when a prior
evaluation has been invalidated by drift---would keep the evidence base current and
directly serve post-market monitoring obligations.

\textbf{A living mapping.} Finally, because the frameworks themselves evolve, the mapping
in this article is a snapshot. We intend to maintain it as a versioned, machine-readable
companion artifact so that the community can extend, correct, and re-time the
method-to-requirement assignments as standards mature and new evaluation methods appear.
